\section{The finite negative-type semimetric}
\label{sec:repair}
\label{sec:metric}

The vectors $P_v$ from \Cref{ker:angles} have the desired macroscopic
distances. Their inner products in \eqref{rep:eq:angle} can still be
negative, but the lemma bounds their negative parts both by $r$ and
by a quantity proportional to within-chart coordinate distance.
We complete the construction by appending a Hilbert correction that
compensates the smaller of these bounds. A correction controlled only
by $r$ could dominate the cost of crossing an individual cell interface.

The following lemma corrects violations of \eqref{rep:eq:angle} while
preserving small distances up to a logarithmic loss. We state it for
arbitrary chartwise $\ell_1$ distances so that the required local and
global bounds are explicit.

\begin{lemma}[Triangle repair]\label{lem:repair}
Let a finite set $V$ be partitioned into charts. In each chart let
$\delta$ be a nonnegative multiple of an $\ell_1$ distance, and set
$\delta(v,w)=\infty$ for points in different charts.
Suppose that $r>0$ and Hilbert vectors $P_v$ satisfy
\begin{equation}\label{rep:eq:hypothesis}
\langle P_v-P_z,P_y-P_z\rangle
\ge-\min\{r,\delta(v,z),\delta(y,z)\}
\qquad(v,y,z\in V).
\end{equation}
Then there are Hilbert vectors $q_v$ such that
\[
d(v,w)=|P_v-P_w|^2+|q_v-q_w|^2
\]
is a negative-type semimetric. They satisfy
\begin{equation}\label{rep:eq:distance}
|q_v-q_w|^2
=8\int_0^r\bigl(1-e^{-\delta(v,w)/l}\bigr)\,dl\le8r.
\end{equation}
If $0<\delta(v,w)\le r$, then
\begin{equation}\label{rep:eq:local}
|q_v-q_w|^2
\le8\delta(v,w)\left(1+\log\frac r{\delta(v,w)}\right).
\end{equation}
If $\delta(v,w)=0$, then $q_v=q_w$.
\end{lemma}

\begin{proof}
We use the convention $e^{-\infty}=0$. Define a Gram matrix by
\begin{equation}\label{rep:eq:gram}
\langle q_v,q_w\rangle=4\int_0^r e^{-\delta(v,w)/l}\,dl.
\end{equation}
This matrix is positive semidefinite. Indeed, in one coordinate the
map $t\mapsto\mathbf1_{[b,t]}$ into $L_2$ has squared distance
$|t-t'|$, where $b$ is a fixed lower bound on the finitely many
coordinate values. Direct sums and scaling show that an $\ell_1$
distance is squared Hilbertian. For finitely many Hilbert vectors
$h_v$, the kernel
\[
e^{-|h_v-h_w|^2/l}
=\E\cos\!\left(\sqrt{2/l}\,\langle G,h_v-h_w\rangle\right)
\]
is positive semidefinite: $G$ is a standard Gaussian in their finite
linear span, and the cosine is the inner product of the corresponding
$(\cos,\sin)$ pairs. This is the Gaussian-kernel criterion of
Schoenberg~\cite{Sch38}, here in its elementary finite form.
The kernels in \eqref{rep:eq:gram} are block diagonal across charts;
integration preserves positive semidefiniteness. Thus the specified
Gram matrix has a Hilbert realization.

The extended distance $\delta$ satisfies the triangle inequality,
including triples meeting different charts. Consequently, for
$a=\delta(v,z)$ and $b=\delta(y,z)$,
\[
\begin{split}
1-e^{-a/l}-e^{-b/l}+e^{-\delta(v,y)/l}
&\ge 1-e^{-a/l}-e^{-b/l}+e^{-(a+b)/l}\\
&=(1-e^{-a/l})(1-e^{-b/l}).
\end{split}
\]
For $t=\min\{r,a,b\}$, this product is at least $(1-e^{-1})^2$
when $0<l<t$. Equation \eqref{rep:eq:gram} therefore gives
\[
\langle q_v-q_z,q_y-q_z\rangle
\ge4(1-e^{-1})^2t\ge t.
\]
Together with \eqref{rep:eq:hypothesis}, this proves
\eqref{rep:eq:angle} for $Z_v=(P_v,q_v)$.

The diagonal entries of \eqref{rep:eq:gram} are $4r$, which proves
\eqref{rep:eq:distance}. Finally, $1-e^{-s}\le\min\{1,s\}$ and
\[
\int_0^r\min\{1,t/l\}\,dl=t\bigl(1+\log(r/t)\bigr)
\qquad(0<t\le r)
\]
give \eqref{rep:eq:local}. The assertion at distance zero follows
directly from \eqref{rep:eq:distance}.
\end{proof}
